\begingroup
\renewcommand{\arraystretch}{1.45}
\fancyhead[L]{\small ARITHMETISCH-GEOMETRISCHES MITTEL.}
\Needspace{28\baselineskip}
\AGMArticle{19}

[Die Beziehungen zwischen den Differentialen der zu untersuchenden Größen sind besonders einfach, wenn $a$ und $b$ unverändert bleiben; dann ergibt sich unmittelbar aus den Bedingungsgleichungen zwischen Gliedern mit gleichem Index, dass der Ausdruck
$\frac{1}{2^n}\frac{1}{b^{\mathrm n}}\surd{\frac{\beta^{\mathrm n}\delta_{\mathrm n}}{\alpha^{\mathrm n}\gamma^{\mathrm n}}}\cdot\mathrm{d}\log\frac{\delta^{\mathrm n}}{\beta^{\mathrm n}}$
seinen Wert nicht ändert, wenn darin der Reihe nach $\alpha,\allowbreak \beta,\allowbreak \gamma,\allowbreak \delta$ statt $\beta,\allowbreak \gamma,\allowbreak \delta,\allowbreak \alpha$ eingesetzt werden. Die beiden so erhaltenen Ausdrücke lassen sich zu einem dem ersten entsprechenden Ausdruck für den Index $n+1$ vereinigen; der Wert dieses Ausdrucks ist also auch vom Index $n$ unabhängig. Durch den Übergang zum Grenzwert und den Vergleich mit den Differentialen der auf verschiedene Weise gebildeten Quotienten ergibt sich
\Needspace{24\baselineskip}
\[
\begin{aligned}\dfrac{1}{k}\mathrm{d}\log\eta
&=\dfrac{1}{a}\surd{\dfrac{\alpha\delta}{\beta\gamma}}\cdot\mathrm{d}\log\dfrac{\delta}{\alpha}
\\[7pt]&=\dfrac{1}{b}\surd{\dfrac{\beta\delta}{\alpha\gamma}}\cdot\mathrm{d}\log\dfrac{\delta}{\beta}
\\[7pt]&=\dfrac{1}{c}\surd{\dfrac{\gamma\delta}{\alpha\beta}}\cdot\mathrm{d}\log\dfrac{\delta}{\gamma}
\\[7pt]&=\dfrac{1}{c}\surd{\dfrac{\delta'}{\beta'}}\cdot\mathrm{d}\log\dfrac{\delta\delta}{\delta'}
\\[7pt]&=\dfrac{1}{a}\surd{\dfrac{\beta\gamma}{\alpha\delta}}\cdot\mathrm{d}\log\dfrac{\gamma}{\beta}
\\[7pt]&=\dfrac{1}{b}\surd{\dfrac{\alpha\gamma}{\beta\delta}}\cdot\mathrm{d}\log\dfrac{\gamma}{\alpha}
\\[7pt]&=\dfrac{1}{c}\surd{\dfrac{\alpha\beta}{\gamma\delta}}\cdot\mathrm{d}\log\dfrac{\alpha}{\beta}
\\[7pt]&=\dfrac{1}{c}\surd{\dfrac{\beta'}{\delta'}}\cdot\mathrm{d}\log\dfrac{\beta'}{\beta\beta}\end{aligned}
\]
\Needspace{16\baselineskip}
eine Beziehung, die also immer gilt, wenn die Indizes von $a,\allowbreak b,\allowbreak c,\allowbreak \alpha,\allowbreak \beta,\allowbreak \gamma,\allowbreak \delta,\allowbreak \eta$ um gleich viele Einheiten erhöht werden. Hieraus folgt der von \textsc{Gauss} aufgezeichnete Satz:]
\Needspace{12\baselineskip}
\[
\begin{aligned}\int\dfrac{\mathrm{d}U}{\surd{(aa\sin U^2+bb\cos U^2)}}
&=\int\dfrac{\mathrm{d}V}{\surd{(aa\cos V^2+bb\sin V^2)}}\\[6pt]&=\dfrac{u}{k}\end{aligned}
\]

\Needspace{22\baselineskip}
\SourceLeaf{400}
[Mithilfe der Gleichung für $\mathrm{d}\log\frac{\beta'}{\beta\beta}$ lassen sich durch ihre wiederholte Anwendung unmittelbar die Reihen für die Differentiale der einzelnen Größen $\alpha,\allowbreak \beta,\allowbreak \gamma,\allowbreak \delta$ aufstellen; für ein negatives $\delta$ entsteht:
\Needspace{22\baselineskip}
\[
\begin{aligned}\dfrac{k}{2}\dfrac{\mathrm{d}\log\dfrac{\alpha}{x}}{\mathrm{d}u}+2c'\sin 2V'
&=\dfrac{k}{2}\dfrac{\mathrm{d}\log\dfrac{\beta}{x}}{\mathrm{d}u}
\\[7pt]&=\dfrac{k}{2}\dfrac{\mathrm{d}\log\dfrac{\gamma}{x}}{\mathrm{d}u}+a\dfrac{\sin V}{\cos U}
\\[7pt]&=\dfrac{k}{2}\dfrac{\mathrm{d}\log\dfrac{\delta}{x}}{\mathrm{d}u}-b\dfrac{\cos V}{\sin U}
\\[7pt]&=c'\sin 2V'+c''\sin 2^2V''+c'''\sin 2^3V'''+.\;.\;.\text{]}\end{aligned}
\]

\Needspace{24\baselineskip}
\AGMArticle{20}

[Die Differentiale zweiter Ordnung, die sich auf die Größe $\eta$ als einzige unabhängige Variable beziehen, lassen sich unmittelbar durch Differentiation der Differentialgleichung erster Ordnung herleiten. Die Bestimmung ihrer Werte lässt sich so darstellen, dass die Ausdrücke von der Form
\Needspace{12\baselineskip}
\[
\dfrac{\mathrm{d}}{\mathrm{d}\log\eta}
\left[\dfrac{\mathrm{d}\log\dfrac{\alpha}{\beta}}{\mathrm{d}\log\eta}\right]
+\tfrac12\dfrac{\mathrm{d}\log\dfrac{\alpha\beta}{\gamma\delta}}{\mathrm{d}\log\eta}
\cdot\dfrac{\mathrm{d}\log\dfrac{\alpha}{\beta}}{\mathrm{d}\log\eta}
\]
für die mit einem beliebigen gemeinsamen Index versehenen $\alpha,\allowbreak \beta,\allowbreak \gamma,\allowbreak \delta$ verschwinden, in welcher Reihenfolge $\alpha,\allowbreak \beta,\allowbreak \gamma,\allowbreak \delta$ auch darin eingesetzt werden. Multipliziert man Zähler und Nenner innerhalb des zweimal zu differenzierenden Logarithmus mit $\beta$, ersetzt $\alpha\beta$ durch $b'\beta'$ und die Ableitungen erster Ordnung durch ihre Werte, so erhält man für zwei aufeinanderfolgende Indizes:
\Needspace{12\baselineskip}
\[
\dfrac{\mathrm{dd}\log\dfrac{\beta'}{\beta\beta}}{(\mathrm{d}\log\eta)^2}
+2\dfrac{a'c'}{kk}\dfrac{\delta'}{\beta'}
-2\dfrac{ac}{kk}\dfrac{\delta}{\beta}-\tfrac12\dfrac{cc}{kk}=0
\]
\Needspace{24\baselineskip}
und hieraus durch wiederholte Anwendung dieser Gleichung und mithilfe der in Artikel 13 für $\mathrm{d}\log\mathrm{M}(a,\allowbreak b)$ gefundenen Reihe die Gleichung
\Needspace{22\baselineskip}
\[
\begin{aligned}\dfrac{\mathrm{d}\log\dfrac{b}{k}}{\mathrm{d}\log y}
&=\dfrac{\mathrm{dd}\log\dfrac{\alpha}{x}}{(\mathrm{d}\log\eta)^2}-\tfrac12\dfrac{ac}{kk}\dfrac{\gamma}{\alpha}
\\[7pt]&=\dfrac{\mathrm{dd}\log\dfrac{\beta}{x}}{(\mathrm{d}\log\eta)^2}+\tfrac12\dfrac{ac}{kk}\dfrac{\delta}{\beta}
\\[7pt]&=\dfrac{\mathrm{dd}\log\dfrac{\gamma}{x}}{(\mathrm{d}\log\eta)^2}-\tfrac12\dfrac{ac}{kk}\dfrac{\alpha}{\gamma}
\\[7pt]&=\dfrac{\mathrm{dd}\log\dfrac{\delta}{x}}{(\mathrm{d}\log\eta)^2}+\tfrac12\dfrac{ac}{kk}\dfrac{\beta}{\delta}\end{aligned}
\]
die gültig bleibt, wenn $a,\allowbreak b,\allowbreak c,\allowbreak \alpha,\allowbreak \beta,\allowbreak \gamma,\allowbreak \delta,\allowbreak x,\allowbreak y,\allowbreak \eta$ mit einem beliebigen gemeinsamen, von null verschiedenen Index versehen werden.

\Needspace{16\baselineskip}
\SourceLeaf{401}
Durch die Verbindung dieses Ergebnisses mit der im vorhergehenden Artikel gefundenen Entwicklung von $\mathrm{d}\log\frac{\alpha}{x}$ ergibt sich die von \textsc{Gauss} gefundene Auswertung des Integrals:]
\Needspace{16\baselineskip}
\[
\begin{aligned}
\int\surd{(aa\sin U^2+bb\cos U^2)}\,\mathrm{d}U
\\[7pt]&=\int\dfrac{aabb\,\mathrm{d}V}{(aa\cos V^2+bb\sin V^2)^{\frac32}}\\
&=\dfrac{u}{k}\{a'a'-2c''c''-4c'''c'''-8c''''c''''-.\;.\}\\
&\quad+c'\sin 2V'-c''\sin 4V''\\&\quad-c'''\sin 8V'''-c''''\sin 16V''''-.\;.\;.\end{aligned}
\]

\clearpage
\endgroup
